% Part: incompleteness
% Chapter: theories-computability
% Section: computably-axiomatizable

\documentclass[../../../include/open-logic-section]{subfiles}

\begin{document}

\olfileid[fr]{inc}{tcp}{cax}

\olsection{Théories axiomatisables}

Une théorie~$\Th{T}$ est dite \emph{!!{axiomatizable}}
si elle admet un ensemble calculable d'axiomes~$A$.
Dire que~$A$ est un ensemble d'axiomes de~$\Th{T}$
signifie que $T=\Setabs{!A}{A \Proves !A}$.
Toute axiomatisation « raisonnable » des entiers
naturels possède cette propriété. En particulier,
toute théorie axiomatisée par un ensemble fini est
!!{axiomatizable}.

\begin{lem}
Toute théorie !!{axiomatizable}~$\Th{T}$ est
!!{computably enumerable}.
\end{lem}

\begin{proof}
Supposons que~$A$ soit un ensemble calculable
d'axiomes de~$\Th{T}$. Pour reconnaître qu'un
!!{sentence}~$!A$ appartient à~$T$, cherchons une
!!{derivation} de~$!A$ à partir de ces axiomes.
Cette recherche aboutit lorsque $!A\in T$.

Autrement dit, $!A$ appartient à~$\Th{T}$ si et
seulement s'il existe une liste finie d'axiomes
$!B_1$, \dots, $!B_k$ de~$A$ et une
!!{derivation}, en logique du premier ordre, de
$(!B_1\land\dots\land !B_k)\lif !A$ (la conjonction vide vaut $\ltrue$).
Nous savons qu'un ensemble défini par une condition
de la forme « il existe \dots\ tel que \dots » est
!!{c.e.} dès lors que la condition placée après
« tel que » est calculable. C'est le cas ici, puisque
l'appartenance à~$A$ et la vérification d'une
!!{derivation} sont calculables.
\end{proof}

\end{document}
